\hypertarget{introduction}{%
\section{Introduction}\label{introduction}}

\hypertarget{problem-statement}{%
\subsection{Problem Statement}\label{problem-statement}}

A confidence-only cascade router like RouteLLM (Ong et al., 2024) reads
a request and decides whether a cheap or expensive model should answer
it. Its router is trained on human preference data, so the only thing it
reads at request time is the request itself. When an edge device answers
that request, it is under continuous CPU and thermal stress. This kind
of stress can allow the present-day neural networks to produce wrong
results (Guo et al., 2017). Even though the softmax output shows the
model is confident, the prediction has already degraded. A router that
only considers confidence cannot identify this specific issue and thus
cannot differentiate an easy input from a hard one. Thus the system
returns a confidently wrong answer.

This is the first failure mode, Failure A, where the routing decision
never sees the operational health of the edge device at the moment of
the inference. BranchyNet (Teerapittayanon et al., 2016) and MSDNet
(Huang et al., 2018) also share the same problem. They both rely on a
confidence or entropy threshold or a computed budget fixed at the design
time. Neither of the systems reads the real-time condition of the
infrastructure. EdgeBERT (Tambe et al., 2021) does react to the state of
the hardware. But it only looks at its own chip's entropy and
voltage-frequency envelope.

A second failure mode, Failure B, sits on top of the first. Present day
cascade systems can run the same fixed enrichment pipeline over every
request. But no outcome ever travels back to it. The system does not
know whether the SLA was met or not, or whether the cloud could have
handled the request better. So the next request is handled in the same
way for similar conditions. They have no pipeline where the decisions
are made differently based on the edge's current state, and no return
path where the outcome of the decision is reported back to the router.

This whole scenario motivates the research question behind this project.
What is gained by replacing static, passive threshold based routing with
a genuinely active orchestration system that can see live data about its
environment and reason over it to take actions and in the end, weigh the
outcomes of those actions?

\hypertarget{aims-and-objectives}{%
\subsection{Aims and Objectives}\label{aims-and-objectives}}
\begin{enumerate}
\def\labelenumi{\arabic{enumi}.}
    \item Survey the landscape of cascade systems, cloud-edge offloading
    and agentic AI architectures from 2016 to 2026. Compare the artifact
    against RouteLLM, the NSGA-II router, Splitwise, GreenServ and
    PROTEUS.
    \item Model the escalation decisions as a Markov Decision Process.
    Here, the state is a fixed 13 slot vector which includes per edge
    resource stress. Its conditional slots are masked out when the
    dynamic Medallion pipeline does not populate them. The actions are
    \emph{route\_to\_edge}, \emph{escalate\_to\_cloud}, or
    \emph{fetch\_extended\_context}.
    \item Implement the dynamic Medallion pipeline, register the metrics
    in a bronze registry, enrich them conditionally into silver and
    produce a masked gold state vector.
    \item Carry out the Edge Context Protocol through self-reports from
    edge and backward loops.
    \item Use an offline domain randomized Gymnasium simulation to train
    PPO policies. The trained frozen policy is then deployed as its own
    services which is accessible via HTTP endpoints.
    \item Provide a dynamic registration and \emph{list\_tools()} method
    of the MCP skill interface.
    \item Performance test against five core baselines and eight
    literature baselines, evaluate on three different traffic
    conditions, conduct a multi-run ablation study and live test in
    multi region cloud environment.
\end{enumerate}

\hypertarget{contributions}{%
\subsection{Contributions}\label{contributions}}

\begin{enumerate}
\def\labelenumi{\arabic{enumi}.}
\item \textbf{Edge as an actor: }The Edge Context Protocol is a new way
of designing edge devices. Instead of just returning predictions, the
device actively reports its own calibration state. This closes Failure
A, a whole class of error that BranchyNet, MSDNet, EdgeBERT and RouteLLM
cannot even detect.
\item \textbf{Medallion Pipeline with Dynamic Components}: A pipeline
that changes what it measures per request, driven by data from the edge.
Which slots of the state vector carry a value varies, while the vector's
own dimension stays fixed at 13.
\item \textbf{SLA remaining and predicted RTT as per-request signals at
routing level}: Every routing decision sees the current request's own
SLA remaining and predicted network RTT, not an aggregate cluster-level
budget. That granularity is missing from Splitwise (cluster/phase-level
scheduling) and every other cascade baseline studied here. It is a soft,
learned preference. The reward function penalizes an SLA violation
directly, but nothing yet blocks a doomed escalation before it happens.
Section~\ref{design-development-and-implementation} states this distinction plainly. It also names a Silver-layer
signal that was computed for that harder guarantee, but is not yet wired
into any decision path.
\item \textbf{MCP Skills Interface at Inference Latency}: A dynamic
per-request tool discovery and registration pattern that lets a
synchronous ML service make use of infrastructure signals well within
the per-request perception budget. The pattern borrows MCP's discovery
and registration shape. It is not an implementation of the Model Context
Protocol wire format, as Section~\ref{design-development-and-implementation} states plainly.
\item \textbf{Two-way agent loop}: Escalation outcomes feed back into
the edge\textquotesingle s own context state within the request
life-cycle itself, not only through a slow offline retraining loop that
doesn\textquotesingle t exist in this design. Together with the dynamic
pipeline above, this addresses Failure B. Ablation Run 4 measures the
loop\textquotesingle s effect directly. Under the retrained policy that
effect is small, which Section~\ref{conclusion} reports rather than claims as a result.
\item \textbf{Policy Orchestration Layer}: This is a governance layer
above the fleet. For every client, it decides whether to reuse,
fine-tune, or train a new policy. That decision is backed by the
task-capacity promotion test (Bossens and Sobey, 2021), and every choice
is logged permanently so it can be checked later.
\end{enumerate}
